\subsection{定义}

命题 1. --- 设 A 是一个交换 Banach 代数。下列条件等价：

\begin{enumerate}
\def\labelenumi{(\roman{enumi})}
\item
  \(\mathsf{X}(\mathsf{A})\) 上的弱拓扑与 Jacobson 拓扑相一致；
\item
  对每个 \(\chi \in \mathsf{X}(\mathrm{A})\) 与每个弱闭子集 \(\mathrm{F}\)（含于 \(\mathsf{X}(\mathrm{A})\)），使得 \(\chi \notin \mathrm{F}\)，存在 \(x \in \mathrm{A}\)，使得 \(\mathcal{G}(x)\) 在 \(\chi\) 处等于 1，在 \(\mathrm{F}\) 上等于 0；
\item
  对 X(A) 的每个弱紧子集 K 与每个弱闭子集 F，若 K ∩ F = ∅，则存在元素 \(x \in A\)，使得 \(\mathcal{G}(x)\) 在 K 上等于 1，在 F 上等于 0。
\end{enumerate}

设 \(M \subset \mathsf{X}(A)\)。说 M 对 Jacobson 拓扑是闭的，意思是：对每个 \(\chi \in \mathsf{X}(A) - \mathsf{M}\)，存在 \(x \in A\)，使得 \(\mathcal{G}(x)\) 在 M 上为零而不在 \(\chi\) 处为零（I, p.~39 的引理 2）。因此条件 (ii) 意味着 \(\mathsf{X}(A)\) 的每个弱闭子集对 Jacobson 拓扑是闭的，这表明 (ii) \(\Longrightarrow\) (i)。此外，(iii) \(\Longrightarrow\) (ii)，因为集合 \(\{\chi\}\) 在 \(\mathsf{X}(A)\) 中是弱紧的。最后，(i) \(\Longrightarrow\) (iii) 由 I, p.~81 的命题 15 的推论得出。

定义 1. --- 设 A 是一个交换 Banach 代数。若它满足命题 1 的等价条件，则称其为正则的。

注. --- 设 \(\tilde{\mathsf{A}}\) 是由 A 添加单位元 \(e\) 而得到的 Banach 代数。命题 1 的条件 (ii) 表明，若 \(\tilde{\mathsf{A}}\) 正则，则 A 正则。设 A 正则，我们来证明 \(\tilde{\mathsf{A}}\) 正则。考察 \(\mathsf{X}(\tilde{\mathsf{A}})\) 的不相交且弱闭（从而弱紧）的子集 F 与 F′，并构造 \(x \in \tilde{\mathsf{A}}\)，使得 \(\mathcal{G}(x)\) 在 F 上为零而在 F′ 上等于 1。设 \(\chi_0 \in \mathsf{X}(\tilde{\mathsf{A}})\) 是 A 上的零特征标。若 \(\chi_0 \notin \mathrm{F}'\)，则由命题 1 的条件 (iii) 以及关于 A 的假设，存在元素 \(x \in \mathsf{A}\)，使得 \(\mathcal{G}(x)\) 在 F 上为零而在 F' 上等于 1。若 \(\chi_0 \in \mathrm{F}'\)，则 \(\chi_0 \notin \mathrm{F}\)；因此存在元素 \(y \in \mathsf{A}\)，使得 \(\mathcal{G}(y)\) 在 F' 上为零而在 F 上等于 1。于是 \(x = e - y\)（属于 \(\tilde{\mathsf{A}}\)）具有所需的性质。

例. --- 我们回到 I, p.~17 第 2 小节的例子。

局部紧空间 X 上在无穷远处趋于 0 的连续复值函数所成的代数（I, p.~17 的例 3）是正则的（参见 I, p.~36, 例 1）。

 \(n\) 次可微且定义在 \([0, 1]\) 上的函数所成的代数（I, p.~18 的例 4）是正则的（参见 I, p.~36, 例 2）。

若 G 是交换局部紧群，\(\mu\) 是 G 上的一个 Haar 测度，则代数 \(L^{1}(G,\mu)\)（I, p.~19 的例 7）是正则的（参见 II, p.~219, 推论 2）。

在圆盘 \(|z| \leqslant 1\) 上连续且在其内部解析的函数所成的代数（I, p.~20 的例 9）不是正则的（参见 I, p.~193, 习题 6）。

命题 2. --- 设 A 是一个交换幺正则 Banach 代数。设 \(n \geqslant 1\) 是一个整数，\((\mathrm{U}_{1}, \ldots, \mathrm{U}_{n})\) 是 \(\mathsf{X}(\mathsf{A})\) 的一个开覆盖。存在 A 中和为 1 的元素 \(x_{1}, \ldots, x_{n}\)，使得 \(\operatorname{Supp}(\mathcal{G}(x_{i})) \subset \mathrm{U}_{i}\)（其中 \(i = 1, \ldots, n\)）。

我们对 n 用归纳法证明该命题。断言对 n = 1 成立。设 \(n \geqslant 2\)，并且断言对 n - 1 已被证明。

存在 \((\mathrm{V}_{1},\ldots,\mathrm{V}_{n})\)（它是 \(\mathsf{X}(\mathsf{A})\) 的开覆盖），使得对每个 i 有 \(\overline{V}_{i}\subset U_{i}\)。由归纳假设，存在元素 \(x,x_{3},\ldots,x_{n}\in A\)，使得 \(x+x_{3}+\cdots+x_{n}=1\)，且 \(\operatorname{Supp}(\mathcal{G}(x))\subset\mathrm{V}_{1}\cup\mathrm{V}_{2}\)，以及对 \(\operatorname{Supp}(\mathcal{G}(x_{i}))\subset\mathrm{V}_{i}\)，有 \(i\geqslant3\)。令 \(\mathrm{K}=\operatorname{Supp}(\mathcal{G}(x))\subset\mathrm{V}_{1}\cup\mathrm{V}_{2}\)。设 \(K_{1}\)（相应地，\(K_{2}\)）为 K 中不属于 \(V_{1}\)（相应地，\(V_{2}\)）的元素组成的集合。那么 \(K_{1}\) 与 \(K_{2}\) 是 K 的不相交紧子集。由于 Banach 代数 A 是正则的，故存在 \(y\in A\)，使得 \(\mathcal{G}(y)=1\) 在 \(K_{1}\) 上成立，且 \(\mathcal{G}(y)=0\) 在 \(K_{2}\) 上成立。于是 \(\mathcal{G}(xy)\) 在 \(\mathsf{X}(\mathsf{A})-\mathsf{K}\) 与 \(K_{2}\) 上为零，因而 \(\operatorname{Supp}\mathcal{G}(xy)\subset\overline{\mathrm{V}}_{2}\subset\mathrm{U}_{2}\)。同理，\(\mathcal{G}(x(1-y))\) 在 \(\mathsf{X}(\mathsf{A})-\mathsf{K}\) 与 \(K_{1}\) 上为零，因而 \(\operatorname{Supp}\mathcal{G}(x(1-y))\subset\overline{\mathrm{V}}_{1}\subset\mathrm{U}_{1}\)。于是元素 \(x_{1}=x(1-y)\)、\(x_{2}=xy\) 与 \(x_{3},\ldots,x_{n}\) 满足命题所述的性质。

推论 1. --- 设 A 是一个交换幺正则 Banach 代数，I 是 A 的理想，\(f: \mathsf{X}(\mathsf{A}) \to \mathbf{C}\) 是连续函数。假设对每个 \(\chi \in \mathsf{X}(\mathsf{A})\)，存在元素 \(y_{\chi} \in \mathrm{I}\)，使得在 \(f = \mathcal{G}(y_{\chi})\) 的一个邻域中有 \(\chi\)。则存在元素 \(y \in \mathrm{I}\)，使得 \(f = \mathcal{G}(y)\)。

由于 X(A) 紧，存在 X(A) 的有限开覆盖 \((\mathrm{U}_{1},\ldots,\mathrm{U}_{n})\)，以及 I 的元素 \(y_{1},\ldots,y_{n}\)，使得在 \(f=\mathcal{G}(y_{i})\) 上有 \(U_{i}\)。由命题 2，存在 A 中和为 1 的元素 \(x_{1},\ldots,x_{n}\)，使得对每个 i 有 \(\operatorname{Supp}(\mathcal{G}(x_{i}))\subset\mathrm{U}_{i}\)。令 \(y=x_{1}y_{1}+\cdots+x_{n}y_{n}\)。这是 I 中具有所需性质的元素。事实上，设 \(\chi\in\mathrm{X}(A)\)。对 \(1\leqslant i\leqslant n\)，我们有 \(\mathcal{G}(x_{i})(\chi)\mathcal{G}(y_{i})(\chi)=\mathcal{G}(x_{i})(\chi)f(\chi)\)，因为若 \(\mathcal{G}(y_{i})(\chi)=f(\chi)\)，则 \(\chi\in U_{i}\)，而若 \(\mathcal{G}(x_{i})(\chi)=0\)，则 \(\chi\notin U_{i}\)。由此得出

\[
\mathcal {G} (y) (\chi) = \sum_ {i = 1} ^ {n} \mathcal {G} (x _ {i}) (\chi) \mathcal {G} (y _ {i}) (\chi) = f (\chi) \sum_ {i = 1} ^ {n} \mathcal {G} (x _ {i}) (\chi) = f (\chi).
\]

推论 2. --- 设 A 是一个交换正则 Banach 代数，I 是 A 的理想，\(f\colon\mathsf{X}'(\mathrm{A})\to\mathbf{C}\) 是连续函数。假设对每个 \(\chi\in\mathsf{X}'(\mathrm{A})\)，存在元素 \(y_{\chi}\in\mathrm{I}\)，使得在 \(f=\mathcal{G}'(y_{\chi})\) 的一个邻域中有 \(\chi\)。则存在元素 \(y\in\mathrm{I}\)，使得 \(f=\mathcal{G}'(y)\)。

设 \(\tilde{\mathsf{A}}\) 是由 \(\mathsf{A}\) 添加单位元而得到的 Banach 代数。那么 \(\tilde{\mathsf{A}}\) 是正则的（注 1），且 \(\mathsf{X}'(\mathsf{A}) = \mathsf{X}(\widetilde{\mathsf{A}})\)；因此只需对 \(\tilde{\mathsf{A}}\) 与理想 I 应用推论 1。

若 I 是交换 Banach 代数的理想，回顾（参见 I, p.~30）：我们用 V(I) 表示核包含 I 的 \(\chi \in \mathsf{X}(\mathrm{A})\) 组成的集合，换言之，即这样的 \(\chi \in \mathsf{X}(\mathrm{A})\) 组成的集合：所有对应于 \(\mathcal{G}(x)\)（\(x \in \mathrm{I}\)）的函数都在其上为零。这是 \(\mathsf{X}(\mathrm{A})\) 的一个对 Jacobson 拓扑为闭的子集。

命题 3. --- 设 A 是一个交换正则 Banach 代数，I 是 A 的理想，K 是 X(A) 的与 V(I) 不相交的紧子集。存在元素 \(x \in I\)，使得 \(\mathcal{G}(x) = 1\) 对 K 中的每个 \(x\) 成立。

这是 I, p.~81 的命题 15 的一个特殊情形，只需注意到 Jacobson 拓扑与 X(A) 上的弱拓扑相一致。
% semantic-fragments: legacy-input-seam
\relax{}
